\chapter{Допуск комплексной симплектической поляризации}
\label{ch:v8-horizontal-phase-complex-symplectic-polarization-admission-gate}

Real-завершение не создаёт независимой обратной стрелки. Проверим теперь
более широкую возможность: допускает ли сам двадцатимерный комплексный
носитель переноса калибровочно-инвариантную невырожденную кососимметрическую
форму.

По полным концевым меткам его представление раскладывается как
\begin{equation}
 \mathcal T_{20}
 \simeq H_+\oplus H_-^{\oplus4}
 \oplus C_+\oplus C_-\oplus S_0^{\oplus4},
 \label{eq:v8-symplectic-transfer-decomposition}
\end{equation}
где
\begin{equation}
 H_+=(\mathbf1,\mathbf2)_{1/2},\quad
 H_-=(\mathbf1,\mathbf2)_{-1/2},\quad
 C_+=(\mathbf3,\mathbf1)_{2/3},\quad
 C_-=(\overline{\mathbf3},\mathbf1)_{-2/3},\quad
 S_0=(\mathbf1,\mathbf1)_0.
 \label{eq:v8-symplectic-transfer-irreps}
\end{equation}
Размерности дают $2+8+3+3+4=20$.

Инвариантная кососимметрическая форма может спаривать только двойственные
изотипические компоненты. Поэтому пространство допустимых форм имеет
точную размерность
\begin{equation}
 \dim\bigl(\Lambda^2\mathcal T_{20}^*\bigr)^G
 =1\cdot4+1\cdot1+\binom42=11.
 \label{eq:v8-symplectic-invariant-form-dimension}
\end{equation}
Это число независимо воспроизведено как нульность точной системы размера
$5600\times190$, полученной из условий
\begin{equation}
 \rho(X)^{\mathsf T}\Omega+\Omega\rho(X)=0,
 \qquad \Omega^{\mathsf T}=-\Omega,
 \label{eq:v8-symplectic-invariance-system}
\end{equation}
для четырнадцати образующих комплексифицированного калибровочного действия.

Существование ненулевых форм не означает существования симплектической
структуры. Слабые двойственные кратности равны $1:4$, поэтому их вклад в
ранг не превосходит $4$. Цветная пара даёт ранг $6$, четыре нейтральных
синглета --- ранг $4$. Следовательно,
\begin{equation}
 \rank\Omega\leq4+6+4=14,
 \qquad
 \dim\ker\Omega\geq20-14=6.
 \label{eq:v8-symplectic-rank-radical-bound}
\end{equation}
Точный общий представитель одиннадцатипараметрического пространства имеет
ранг $14$, так что оценка достигается и является строгой.

Поляризация также не единственна. Например, можно спарить $H_+$ с первой
или со второй из четырёх копий $H_-$. Обе формы калибровочно-инвариантны,
кососимметричны и имеют ранг $14$, но различны:
\begin{equation}
 \Omega_1\neq\Omega_2,
 \qquad
 \rank\Omega_1=\rank\Omega_2=14.
 \label{eq:v8-symplectic-two-polarizations}
\end{equation}

Есть и второй независимый запрет. Для одного бозонного поля с коммутирующими
координатами любая такая форма даёт
\begin{equation}
 \Phi^{\mathsf T}\Omega\Phi=0.
 \label{eq:v8-symplectic-bosonic-self-contraction}
\end{equation}
Ненулевой голоморфный скаляр потребовал бы второго независимого поля либо
некоммутирующих или грассмановых координат.

Минимальная невырожденная достройка обязана уравнять слабые кратности. Для
этого недостаёт трёх копий $H_+$, то есть
\begin{equation}
 \Delta\dim_{\mathbb C}=3\cdot2=6,
 \qquad
 \dim_{\mathbb C}\mathcal T_{\rm symp}=26.
 \label{eq:v8-symplectic-minimal-completion}
\end{equation}
Эти шесть комплексных направлений не содержатся в нынешнем концевом
носителе и требуют нового допуска.

\begin{theorem}[Вырожденное семейство инвариантных форм]
На полном двадцатимерном носителе существует одиннадцатимерное пространство
калибровочно-инвариантных кососимметрических форм. Ни одна из них не является
невырожденной: максимальный ранг равен $14$, а минимальный радикал имеет
размерность $6$. Калибровочная симметрия не выбирает единственную форму.
\end{theorem}

\begin{nogocriterion}[No-go текущей симплектической поляризации]
Нельзя объявлять одну из одиннадцати форм канонической и нельзя использовать
$\Phi^{\mathsf T}\Omega\Phi$ как фазовую массу: эта свёртка тождественно
нулевая. Невырожденная симплектическая архитектура требует как минимум шесть
новых комплексных слабых двойственных направлений и отдельного второго
поля либо отображения момента.
\end{nogocriterion}

\begin{protocol}
\begin{enumerate}
 \item комплексная размерность носителя: \textbf{$20$};
 \item инвариантных кососимметрических форм: \textbf{$11$};
 \item максимальный ранг формы: \textbf{$14$};
 \item минимальная размерность радикала: \textbf{$6$};
 \item две различные максимальные формы: \textbf{предъявлены};
 \item каноническая поляризация: \textbf{не выбрана};
 \item бозонная самосвёртка: \textbf{$0$};
 \item недостающие двойственные направления: \textbf{$6$ комплексных};
 \item минимальная симплектическая достройка: \textbf{$26$ комплексных};
 \item подъём горизонтальной моды на текущем носителе: \textbf{нет};
 \item следующий гейт: \textbf{концевой допуск минимальной симплектической
 достройки}.
\end{enumerate}
\end{protocol}

Машинный сертификат сохранён в
\path{s2t_v8_horizontal_phase_complex_symplectic_polarization_admission_gate_results.json}.

\begin{thebibliography}{9}
\bibitem{v8-symplectic-quiver-moment-map}
M.~Harada, G.~Wilkin,
\emph{Morse theory of the moment map for representations of quivers},
Geom. Dedicata 150 (2011), 307--353; arXiv:0807.4734.
\bibitem{v8-symplectic-preprojective}
W.~Crawley-Boevey,
\emph{Poisson structures on moduli spaces of representations},
J. Algebra 325 (2011), 205--215; arXiv:math/0502301.
\end{thebibliography}